% Section .............................................................................................................
\section{Entanglement, Structural Realism, and Conclusion}
\label{section:structural}

% Subsection ..........................................................................................................
\subsection{Entanglement as Irreducibly Relational Potentiality}

The potentiality reading connects naturally to ontic structural realism (OSR) \citep{ladyman_2007_everything, french_2014_structure}. OSR proposes that what is real
are structures---patterns of relations---rather than objects with intrinsic natures, motivated in part by the permutation symmetry of
identical particles \citep{french_1988_indiscernibles} and by the non-locality revealed by Bell inequality violations \cite{bell_1964_epr, aspect_1982_bell, hensen_2015_loophole}.

\medskip

Consider the spin-singlet state of two electrons,
\begin{equation}
  \ket{\Psi^-} = \frac{1}{\sqrt{2}}\bigl(\ket{\uparrow}_A\ket{\downarrow}_B - \ket{\downarrow}_A\ket{\uparrow}_B\bigr).
  \label{eq:bell}
\end{equation}

This state cannot be written as a product $\ket{\phi}_A\otimes\ket{\chi}_B$. Neither particle $A$ nor $B$ has a definite spin potentiality on its own. On the Aristotelian reading, the potentialities are properties not of $A$ and $B$ separately but of the \emph{pair} $AB$ jointly---an irreducibly relational structure.

\medskip

Bell's theorem \citep{bell_1964_epr} shows that no local hidden-variable theory can reproduce the quantum correlations. The CHSH inequality
\begin{equation}
  \bigl|E(a,b) - E(a,b') + E(a',b) + E(a',b')\bigr| \leq 2
  \label{eq:chsh}
\end{equation}
is violated by quantum mechanics up to $2\sqrt{2}$ (the Tsirelson bound), confirmed experimentally \citep{aspect_1982_bell, hensen_2015_loophole}. Entanglement is not a correlation between independently real objects. It is a structure of irreducibly joint potentiality with no local decomposition---precisely the structural realist's claim that relations are ontologically prior to relata.

% Subsection ..........................................................................................................
\subsection{Potentiality Realism vs.\ Ontic Structural Realism}

The question arises whether the Aristotelian potentiality reading simply \emph{is} a form of OSR. We suggest it occupies a distinct position. OSR in its strongest form eliminates objects entirely in favour of structures \citep{french_1988_indiscernibles}. The Aristotelian reading preserves objects---the electron is a genuine entity---but denies that objects are fully actual with respect to all their properties prior to measurement.

\medskip

This is a form of \emph{potentiality realism}: realism about both objects and their potentialities, with actuality restricted to what has been, or is being, measured. Objects are real; their potentialities are real; their actualities are contingent on measurement interactions.

\medskip

\TODO{Engage with Esfeld and Deckert's minimalist ontology \citep{esfeld_2017_minimalist}: points in a structure with no intrinsic properties. Argue potentiality realism is distinct: it accepts intrinsic properties \emph{in potentia}, not \emph{in actu}.}

% Subsection ..........................................................................................................
\subsection{Conclusion}

The classical calculation of the hydrogen atom's radiative collapse is more than a historical curiosity. It is a precise, mathematically impeccable derivation---from Newton, Coulomb, and Larmor, all independently confirmed in their domains---that produces an empirically false result. And the reason it is false is not a flaw in the equations. It is a flaw in the ontology they presuppose.

\medskip

We have argued for three connected claims. First, every physical model carries an ontological layer analytically distinct from its predictive-mathematical content; diagnosing which layer has failed is a non-trivial philosophical task that physics alone cannot perform. Second, Bohr's 1913 resolution was an explicit ontological stipulation---he declared the classical ontology inapplicable without fully replacing it---and the residue of this incompleteness persists in the renormalisation programme of QED. Third, Aristotle's potentiality/actuality distinction, recovered by Heisenberg in his 1958 essay, provides a more adequate starting ontology for quantum mechanics than either classical particle realism or strict instrumentalism. Finally we discuss entanglement, which read through this lens, is a structure of irreducibly joint potentiality that connects naturally to ontic structural realism while preserving a genuine role for objects.

\medskip

The lesson is general. Scientific progress is not only the accumulation of more accurate predictions. It is the successive replacement of ontologies---each adequate to a domain, each carrying implicit metaphysical commitments, each eventually confronted with phenomena that force not just new equations but a new picture of what there is.

\medskip

Sixteen picoseconds. That is how long a classical hydrogen atom would survive. That it does not survive that way, and that we needed a century of physics and philosophy to understand why, says something important about the relationship between mathematics, experiment, and the picture of the world that any physical theory always---whether acknowledged or not---carries with it.
